\chapter{ВАЛИДАЦИЯ, ИМПАКТ И БИЗНЕС-ОЦЕНКА}

% ─────────────────────────────────────────────────────────────────
\section{Протокол валидации}
% ─────────────────────────────────────────────────────────────────

\subsection{Автоматизированная оценка}

Качество выходных артефактов оценивается автоматически
в~трёх измерениях:

\begin{enumerate}
  \item \textbf{Ragas-метрики} (глава~2): достоверность относительно
    источника (faithfulness), релевантность ответа (answer~relevancy),
    точность контекста (context~precision);
  \item \textbf{Покрытие уровней таксономии Блума}: доля уровней таксономии Блума,
    покрытых в~наборе учебных целей курса; целевой показатель~--- $\geq 70\%$;
  \item \textbf{Структурное качество}: полнота (наличие всех обязательных
    элементов методологии), согласованность учебных целей с~содержанием,
    BLEU/ROUGE относительно эталонов B1.
\end{enumerate}

\subsection{Протокол экспертной оценки}

Независимая верификация проводится 5--8~экспертами-методистами
(отбор: опыт разработки eLearning-курсов~$\geq$~3~лет,
опыт оценки учебных программ).

Каждому эксперту предъявляются четыре обезличенных варианта структуры
одного курса (B1--B4) в~случайном порядке. Оценка проводится
по~рубрике из~5~критериев (шкала 1--5):

\begin{itemize}
  \item соответствие целевой аудитории;
  \item логика и~последовательность модулей;
  \item качество учебных целей (SMART);
  \item покрытие исходного материала;
  \item практическая применимость.
\end{itemize}

Итоговая оценка --- среднее по~5~критериям.
Слепой дизайн исключает предвзятость в~пользу черновика,
подготовленного системой.
Полные таблицы текущих оценочных артефактов приведены в~приложении~И.

\subsection{Экспертная и пользовательская оценка}

Помимо автоматизированных метрик подготовлены два контура
качественной проверки.

Первый контур --- \textbf{слепая экспертная оценка}. Для~него
подготовлена рубрика из~5 критериев и~шаблон таблицы сравнения
конфигураций B1--B4. Этот контур должен быть проведён после стабилизации
мультиагентной конфигурации B4, чтобы экспертам предъявлялись уже
не~прототипные, а воспроизводимые результаты.

Второй контур --- \textbf{пользовательский протокол пилота}. Он
предполагает измерение времени до первого черновика, субъективного удобства
(SUS~\cite{brooke_sus_1996}) и~готовности рекомендовать решение
(NPS). Дополнительно в~следующую волну пилотов закладываются
опросы на~самоэффективность (task-specific self-efficacy)
и~ощущаемую компетентность (perceived competence), выполненные
в~формате retrospective pre-post, а~также отложенный опрос
на~перенос обучения в~практику (transfer of learning)
и~достижение цели (goal attainment) через 2--4~недели после использования
системы
~\cite{bandura_self_efficacy_2006, pcs_sdt_2026,
kirkpatrick_model_2026, shilts_retrospective_2008,
shogren_gas_2021}. На~момент текущей редакции работы пилотный сценарий
и~вопросы подготовлены, однако полный цикл пользовательской оценки
ещё не~завершён. Это разграничение принципиально: диплом фиксирует уже
доступные технические результаты и~отдельно описывает инструменты
для~следующей волны validation.

% ─────────────────────────────────────────────────────────────────
\section{Сквозная верификация}
% ─────────────────────────────────────────────────────────────────

\subsection{Сценарий~1: Онбординг-курс (ADDIE)}

\textbf{Входной документ:} продуктовое описание нового SaaS-инструмента,
8~страниц.

\textbf{Ожидаемый результат:} 6-модульный онбординг-курс с~учебными
целями для~новых сотрудников; время подготовки первого методического
черновика менее 30~минут.

\textbf{Ход выполнения:}
\begin{enumerate}
  \item Пользователь загружает PDF в~систему; запускается контур обработки документов.
  \item После индексации инициируется CourseStructureAgent
    с~методологией ADDIE.
  \item После~генерации структуры система приостанавливается для~ручного ревью;
    пользователь одобряет структуру.
  \item ContentCreatorAgent генерирует тексты уроков; FactCheckAgent
    верифицирует утверждения.
  \item Verifier формирует quality seal, provenance и snapshot
    дипломных метрик H1--H3.
  \item После финального human-review курс может быть опубликован
    в~LMS Adapstory.
\end{enumerate}

На~текущем этапе количественно измерен не~весь сквозной конвейер B4,
а~его наиболее воспроизводимый базовый слой B2. Для~измеренного корпуса
из~19 документов конфигурация B2 показывает среднее время подготовки структуры
4.37~с, среднее покрытие уровней таксономии Блума 61.4\,\%, долю курсов
с~покрытием $\geq 70\,\%$ на~уровне 21.1\,\% и~структурную валидность
100\,\%. Эти цифры ещё не~являются итоговыми показателями сценария~1,
но подтверждают, что минимальный слой «документ -> структура курса» уже
работоспособен и~может служить отправной точкой для~улучшения B4.

\subsection{Сценарий~2: Нормативный курс с~факт-чекингом}

\textbf{Входной документ:} корпоративный регламент информационной
безопасности, 22~страницы.

\textbf{Специфика:} документ содержит точные числовые требования
(сроки, пороги, санкции), критичные для~корректного нормативного обучения.

Проверяется, что FactCheckAgent выявляет галлюцинации при~конкретизации
числовых значений. Для~этого уже подготовлен синтетический датасет
из~100 утверждений (\texttt{statements.json}), пригодный для~расчёта
precision / recall / F1. Финальные численные результаты будут
получены после стабилизации B4 и~подключения полного контура факт-чекинга
к~основному оценочному прогону.

\subsection[Сценарий 3: Автоактуализация]{Сценарий~3: Авто-актуализация при~обновлении источника}

\textbf{Условие:} в~исходный документ сценария~1 внесены изменения
(обновлена функциональность продукта, 3~новых раздела).

Проверяется корректность выявления устаревших уроков и~качество
перегенерации затронутых разделов без~перегенерации всего курса.
На~момент подготовки работы сценарий сформулирован как~отдельный
validation-case, однако численная метрика ещё не~зафиксирована.
Это один из~следующих экспериментов, поскольку именно он проверяет
прикладную ценность системы для~B2B-контента, который быстро устаревает.

% ─────────────────────────────────────────────────────────────────
\section{Текущие результаты и проверка гипотез}
% ─────────────────────────────────────────────────────────────────

\subsection{Сводка фактически подтверждённых результатов}

\begin{table}[h!]
  \centering
  \small
  \caption{Текущий набор подтверждённых результатов исследования}
  \label{tab:current-results}
  \begin{tabular}{|p{3.9cm}|p{2.7cm}|p{6.2cm}|}
  \hline
  \textbf{Контур валидации}
    & \textbf{Текущее значение}
    & \textbf{Что это означает} \\
  \hline
  Customer Development
    & 20 B2B-интервью
    & проблема подтверждена первичными данными рынка \\
  \hline
  Ключевой pain
    & 70\,\%
    & большинство респондентов называют длительность разработки курса
      главным ограничением текущего процесса \\
  \hline
  Использование подрядчиков / внешних методистов
    & 50\,\%
    & рынок уже платит за~закрытие этой задачи внешним ресурсом \\
  \hline
  Готовность рассмотреть ИИ-инструмент
    & 60\,\%
    & есть ранний сигнал интереса к~решению, но ещё не завершён
      полноценный этап проверки соответствия решения проблеме \\
  \hline
  Базовая конфигурация B2: время подготовки структуры
    & 4.37~с
    & одношаговая конфигурация очень быстра, но не решает задачу качества
      и~объяснимости \\
  \hline
  Базовая конфигурация B2: среднее покрытие уровней Блума
    & 61.4\,\%
    & текущая конфигурация не достигает целевого уровня $\geq 70\,\%$ \\
  \hline
  Базовая конфигурация B2: доля курсов с~Блумом $\geq 70\,\%$
    & 21.1\,\%
    & именно этот разрыв обосновывает необходимость агентного усиления \\
  \hline
  Базовая конфигурация B2: структурная валидность
    & 100\,\%
    & минимальный конвейер уже стабильно выдаёт формально корректную
      структуру курса \\
  \hline
  Оценочные наборы
    & 30 вопросно-ответных пар, 100 утверждений
    & база для~следующей волны RAG-оценки и оценки факт-чекинга уже подготовлена \\
  \hline
  Сценарная unit-экономика
    & 75~руб./курс
    & инфраструктурная себестоимость выглядит низкой даже до~завершения
      полного оценочного цикла B4 \\
  \hline
  Отношение к~подрядной разработке
    & 0.05\,\%
    & расчётная себестоимость существенно ниже стоимости внешней ручной
      разработки курса \\
  \hline
  LTV / CAC, срок окупаемости, безубыточность
    & 10.0; 2~мес.; 14 клиентов
    & сценарная модель выглядит экономически осмысленной, но требует
      подтверждения реальными продажами \\
  \hline
\end{tabular}
\end{table}

Сравнение результатов B2 с~целевыми показателями визуализировано на рисунке~\ref{fig:b2-vs-target}.

% Сравнение результатов B2 с целевыми показателями
\begin{figure}[htbp]
  \centering
  \begin{tikzpicture}
    \begin{axis}[
      ybar,
      width=0.9\textwidth,
      height=0.55\textwidth,
      bar width=18pt,
      ylabel={Значение, \%},
      ymin=0, ymax=120,
      xtick=data,
      symbolic x coords={Покрытие Блума, Доля курсов $\geq$70\%, Структурная валидность},
      x tick label style={align=center, font=\small},
      legend style={at={(0.5,-0.2)}, anchor=north, legend columns=2},
      nodes near coords,
      nodes near coords align={vertical},
      every node near coord/.append style={font=\footnotesize},
      enlarge x limits=0.25,
      grid=major,
      ymajorgrids=true,
      xmajorgrids=false,
    ]
    % B2 результаты
    \addplot[fill=blue!60, draw=blue!80] coordinates {
      (Покрытие Блума, 61.5)
      (Доля курсов $\geq$70\%, 18.8)
      (Структурная валидность, 100)
    };
    % Целевые показатели
    \addplot[fill=gray!30, draw=gray!70, postaction={pattern=north east lines, pattern color=gray!60}] coordinates {
      (Покрытие Блума, 70)
      (Доля курсов $\geq$70\%, 80)
      (Структурная валидность, 100)
    };
    \legend{Результат B2, Целевой показатель}
    \end{axis}
    % Аннотация разрывов
    \node[font=\footnotesize, text=red!70!black, align=center] at (2.1, 5.2) {$\Delta$\,=\,--8.5\,п.п.};
    \node[font=\footnotesize, text=red!70!black, align=center] at (4.9, 3.5) {$\Delta$\,=\,--61.2\,п.п.};
  \end{tikzpicture}
  \caption{Сравнение результатов конфигурации B2 с целевыми показателями}
  \label{fig:b2-vs-target}
\end{figure}


Отдельно от чисел B2 в платформе уже реализован машинно читаемый
контур B4-доказательств: result/status payload содержит
\texttt{citations}, \texttt{quality\_seal}, \texttt{block\_provenance},
\texttt{cost\_quality} и \texttt{thesis\_validation}. Эти поля
покрыты фронтовым E2E-сценарием отображения дипломных метрик и
используются скриптом \texttt{run\_ai\_methodist\_product\_validation.py}
для формирования gate-report. Поэтому текущее ограничение H2 состоит
не в отсутствии платформенного механизма source-grounding, а в том,
что paired B2/B4-прогон на одном корпусе ещё должен быть завершён
как исследовательский эксперимент.

\subsection{Проверка гипотез H1--H3}

\textbf{H1 (продуктовая).}
«Мультиагентный конвейер должен готовить первый методический
черновик курса не дольше чем за~30~минут
при~сохранении покрытия уровней таксономии Блума $\geq 70\%$.»

На~момент текущего черновика гипотеза получает
\textbf{предварительную поддержку}, но не может считаться полностью
подтверждённой. С одной стороны, проблема времени разработки надёжно
валидирована через 20 интервью, а текущая конфигурация B2 показывает, что
генерация структуры курса уже занимает секунды, а не часы. С другой
стороны, полноценный сценарий B4 с~ручным ревью и~итоговым
измерением времени до первого черновика на~реальном пилоте ещё не стабилизирован.
\textbf{Гипотеза H1 предварительно поддержана и требует финального пилотного прогона.}

\textbf{H2 (техническая).}
«Мультиагентная архитектура обеспечивает более высокое качество структуры
по~сравнению с~одношаговой базовой конфигурацией на~той же LLM.»

Текущее состояние данных позволяет утверждать лишь то, что
одношаговая конфигурация B2 недостаточна по~покрытию уровней Блума
и~совсем не обеспечивает явной ссылочной привязки к~источнику.
Это создаёт содержательное основание для
мультиагентной архитектуры, но ещё не даёт статистически завершённого
сравнения B2 и~B4. До завершения стабильного оценочного прогона гипотеза H2
должна трактоваться как \textbf{не полностью проверенная}.
\textbf{Гипотеза H2 концептуально обоснована, но эмпирически пока не завершена.}

\textbf{H3 (бизнесовая).}
«Себестоимость подготовки курса $\leq 1\%$ от~стоимости ручной разработки
у~внешнего подрядчика.»

В~сценарной финансовой модели себестоимость подготовки курса составляет
75~руб., что эквивалентно примерно 0.05\,\% от~ориентировочной стоимости
ручной разработки курса у~внешнего подрядчика.
Такой результат не доказывает рынок сам по~себе, но показывает, что
архитектура не разрушает экономику продукта уже на~раннем этапе.
\textbf{Гипотеза H3 предварительно поддержана сценарной unit-экономикой.}

\subsection{Производительность}

Подтверждённые операционные показатели:
\begin{itemize}
  \item среднее время конфигурации B2 на~измеренном корпусе из~19 документов ---
    4.37~с на~генерацию структуры курса;
  \item среднее число модулей --- 3, среднее число уроков --- 6.2;
  \item доля структурно валидных результатов --- 100\,\%.
\end{itemize}

Нодовая телеметрия p50 / p95, пропускная способность и~загрузка GPU включены
в~план следующей итерации наблюдаемости. Они будут сниматься после
стабилизации B4 через LangSmith / системные метрики рантайма.

\subsection{Точность факт-чекинга}

Синтетический датасет из~100 утверждений (50 корректных,
50 намеренно искажённых) позволяет воспроизводимо рассчитать
precision / recall / F1. Численные значения не включены
в~текущую редакцию, поскольку полный B4-контур факт-чекинга
ещё не завершил стабильный прогон.

% ─────────────────────────────────────────────────────────────────
\section{Ограничения исследования и угрозы валидности}
\label{sec:limitations}
% ─────────────────────────────────────────────────────────────────

Ниже перечислены ограничения исследования и~угрозы
внутренней и~внешней валидности~\cite{wohlin_experimentation_2012},
существенные для~интерпретации результатов.

\subsection{Ограничения тестового датасета}

\textbf{Ограниченный объём измеренного корпуса~($n = 19$).}
Ядровой корпус содержит 19~документов, что ограничивает
статистическую мощность сравнений.
При~размере эффекта $d = 0.5$ (средний) и~$\alpha = 0.05$
мощность критерия Вилкоксона для~такой выборки остаётся ограниченной
и~не~гарантирует надёжного сравнения тонких различий между конфигурациями.
Следовательно, результаты следует рассматривать как~предварительные
и~требующие репликации на~расширенном корпусе ($n \geq 30$).

\textbf{Жанровая и языковая фокусировка измеренного ядра.}
Ядровой корпус сознательно собран как смешанный срез:
русскоязычные продуктово-регламентные документы и англоязычные учебные
программы MIT~OCW. Это повышает прикладную релевантность для
B2B- и образовательных сценариев, но ограничивает прямой перенос
количественных результатов на другие отрасли, регуляторные среды
и типы исходных документов. USWDS, GOV.UK
и~Saylor Academy вынесены в~зарубежный сравнительный контур
и~в~текущей версии используются только качественно, без полного
количественного прогона.

\subsection{Ограничения эталонного корпуса}

На~момент текущей редакции полностью завершённым является не~экспертный
эталонный корпус, а~корпус входных документов и~оценочные наборы
для~RAG и~факт-чекинга. Экспертная разметка и~расчёт согласия
между экспертами подготовлены как~следующий шаг, но ещё не могут считаться
окончательно зафиксированными результатами. Это ограничивает силу
выводов о~педагогическом превосходстве B4 по~отношению к~другим
конфигурациям.

\subsection{Ограничения пользовательской валидации}

Пользовательский пилотный контур спроектирован, но ещё не завершён как
формальный исследовательский этап. Поэтому текущая работа пока не даёт
численно подтверждённых SUS / NPS и~не может делать сильные выводы
о~поведенческом удержании, повторном использовании или~организационном
масштабировании. Следующий шаг --- проведение решенческих интервью,
этапа писем о~намерениях и~пилотов на~реальных документах заказчиков.

\subsection{Ограничения технологического стека}

\textbf{Зависимость от~одной LLM.}
Все~эксперименты проведены на~одной локальной модели семейства Qwen~2.5.
Результаты могут не~воспроизводиться на~других моделях
(LLaMA~3.3, Mistral~Large, GPT-4o), что ограничивает
обобщаемость вывода об~эффективности агентной архитектуры.

\textbf{Ограниченный набор методологий.}
На~момент валидации реализована одна методология (ADDIE).
Сравнение с~SAM и~Backward~Design отложено до~следующей итерации
и~не~включено в~основные результаты работы.

\textbf{Авто-классификация уровней таксономии Блума.}
Уровни таксономии Блума определялись автоматически с~помощью
GPT-4o в~роли~judge-модели~\cite{mt_bench_2023}.
Надёжность такой классификации на~русскоязычном
наборе учебных целей
не~валидирована независимо; это является угрозой валидности
измерения покрытия уровней таксономии Блума.

\subsection{Угрозы внешней валидности}

Система разработана и~количественно апробирована на~русскоязычном
текстовом контенте продуктового, нормативного и~образовательного
характера; зарубежный контур пока использован только качественно.
Перенос результатов на~мультимодальный контент, закрытые внутренние
документы компаний, другие отрасли и~другой регуляторный контекст
(SCORM, xAPI) требует дополнительных исследований.

\subsection{Направления для устранения ограничений}

Следующие шаги:
\begin{enumerate}
  \item расширение российского ядра до~30--50~документов
    из~5+~различных сценариев;
  \item добавление второго эксперта и~расчёт Cohen's~$\kappa$
    для~оценки надёжности разметки уровней таксономии Блума;
  \item кросс-модельная оценка: воспроизведение экспериментов
    на~LLaMA~3.3~70B и~GPT-4o для~проверки обобщаемости;
  \item отдельный количественный прогон на~зарубежном сравнительном
    контуре после стабилизации российского baseline;
  \item расширенный пилот ($n \geq 20$) с~продольными
    метриками удержания контента.
\end{enumerate}

% ─────────────────────────────────────────────────────────────────
\section{Бизнес-оценка и масштабирование}
% ─────────────────────────────────────────────────────────────────

\subsection{Unit-экономика}

Себестоимость одного курса в текущем MVP складывается из~LLM-времени,
retrieval/KG-вызовов, хранения состояния и~опциональных API-вызовов.
Сценарный расчёт собран в~артефакте \texttt{unit\_economics\_results.json}.

\begin{itemize}
  \item GPU-время~--- 75~руб./курс;
  \item OpenAI API (опционально)~--- 0 в~базовой конфигурации;
  \item \textbf{Итоговая себестоимость}~--- $\approx$~75~руб./курс.
\end{itemize}

Соотношение себестоимости ИИ-генерации и~подрядной разработки показано на рисунке~\ref{fig:unit-economics}.

\begin{figure}[htbp]
  \centering
  \begin{tikzpicture}
    \begin{axis}[
      xbar,
      width=0.78\textwidth,
      height=4.5cm,
      bar width=18pt,
      xlabel={Стоимость, руб. (логарифмическая шкала)},
      xmode=log,
      xmin=10, xmax=500000,
      log ticks with fixed point,
      ytick=data,
      symbolic y coords={{ИИ-генерация курса},{Подрядная разработка}},
      y tick label style={font=\small},
      nodes near coords,
      nodes near coords align={horizontal},
      every node near coord/.append style={font=\footnotesize},
      xmajorgrids=true,
      grid style={dashed, gray!40},
      enlarge y limits=0.5,
    ]
    \addplot[fill=teal!60, draw=teal!80] coordinates {
      (75,{ИИ-генерация курса})
    };
    \addplot[fill=red!40, draw=red!60] coordinates {
      (150000,{Подрядная разработка})
    };
      \node[font=\small, text=teal!60!black, anchor=west]
        at (rel axis cs:0.56,0.78) {$\times$2000 разница};
    \end{axis}
  \end{tikzpicture}
  \caption{Unit-экономика ИИ-методиста: себестоимость vs рыночная стоимость}
  \label{fig:unit-economics}
\end{figure}


Метрики сценарной подписочной модели:

\begin{itemize}
  \item CAC (стоимость привлечения клиента): около 30\,000~руб.;
  \item LTV (пожизненная ценность клиента): около 300\,000~руб.;
  \item LTV / CAC: 10.0;
  \item срок окупаемости: 2~мес.
\end{itemize}

\subsection{Lean Canvas}

Lean~Canvas~\cite{maurya_lean_canvas_2012} ИИ-методиста приведён
в~приложении~К. Основные блоки:

\textbf{Проблема}: разработка корпоративного курса занимает много
времени и~часто выносится на~подрядчиков; контент устаревает
быстрее, чем его успевают перерабатывать.

\textbf{Решение}: полуавтоматическая подготовка структуры и~черновых
материалов курса по~загруженному документу; проактивные подсказки
и~ручное ревью в~контуре; механизм последующей авто-актуализации.

\textbf{Целевая аудитория}: B2B-клиенты платформы Adapstory~---
HR-отделы и~учебные центры с~>~50~сотрудников.

\textbf{Труднокопируемое преимущество}: методологическая вариативность (ADDIE/SAM/BD)
и~полный суверенитет данных за~счёт локального LLM-контура не~воспроизводятся
Coursebox~AI и~iSpring~AI.

\subsection{Финансовая модель (прогноз на 3 года)}

Детальные расчёты приведены в~приложении~К.
Итоговые показатели:

\begin{itemize}
  \item \textbf{NPV}: 1\,807~тыс.\ руб.;
  \item \textbf{IRR}: 32.2\,\%;
  \item \textbf{ROI} за~3~года: 42.2\,\%;
  \item \textbf{Точка безубыточности}: 14~клиентов
    / 32~мес.\ с~момента коммерческого запуска
\end{itemize}

Динамика выручки и~расходов представлена на рисунке~\ref{fig:financial-forecast}.

% Финансовый прогноз на 3 года
\begin{figure}[htbp]
  \centering
  \begin{tikzpicture}
    \begin{axis}[
      width=0.9\textwidth,
      height=0.55\textwidth,
      xlabel={Месяц},
      ylabel={тыс. руб./мес.},
      xmin=0, xmax=36,
      ymin=0, ymax=450,
      grid=major,
      grid style={dashed, gray!30},
      legend style={at={(0.3,0.95)}, anchor=north, legend columns=1},
      xtick={0,6,12,18,24,30,36},
    ]
    % Выручка
    \addplot[blue, thick, mark=*, mark size=2pt] coordinates {
      (0, 0)
      (3, 15)
      (6, 45)
      (9, 75)
      (12, 120)
      (15, 150)
      (18, 180)
      (21, 210)
      (24, 240)
      (27, 270)
      (30, 300)
      (33, 330)
      (36, 360)
    };
    \addlegendentry{Выручка}
    % Расходы
    \addplot[red!70!black, thick, dashed, mark=square*, mark size=2pt] coordinates {
      (0, 200)
      (6, 200)
      (12, 200)
      (18, 200)
      (24, 200)
      (30, 200)
      (36, 200)
    };
    \addlegendentry{Расходы}
    % Точка безубыточности — вертикальная линия
    \draw[gray, densely dashed, thick] (axis cs:32,0) -- (axis cs:32,420);
    \node[font=\footnotesize, align=center, anchor=south west] at (axis cs:32.5,350)
      {Точка\\безубыточности\\(32 мес.)};
    % Точка пересечения
    \addplot[only marks, mark=o, mark size=5pt, thick, black] coordinates {(32, 200)};
    \end{axis}
  \end{tikzpicture}
  \caption{Финансовый прогноз на 3 года (сценарная модель)}
  \label{fig:financial-forecast}
\end{figure}


\subsection{Потенциал масштабирования}

После~валидации в~российском B2B-сегменте выделены
смежные рынки:

\begin{itemize}
  \item \textbf{Международный EdTech}: потенциальный TAM $\approx$~\$51~млрд
    к~2030~году~\cite{edtech_market_report_2024};
  \item \textbf{Корпоративное обучение (CLO-сегмент)}: интеграция
    в~Microsoft~365~+~SharePoint как источник документов;
  \item \textbf{Государственное образование}: адаптация под~ФГОС-регламент
    при~поддержке профильных ведомств.
\end{itemize}

\subsection{PESTLE-анализ при масштабировании}

\begin{itemize}
  \item \textbf{P (Political)}: регуляторные требования к~использованию
    ИИ в~государственном образовании; ограничения экспорта технологий;
  \item \textbf{E (Economic)}: зависимость от~стоимости GPU-аренды;
    волатильность курсов валют при~API-зависимости;
  \item \textbf{S (Social)}: опасения преподавателей и~методистов
    относительно снижения роли человека; требования к~прозрачности
    использования ИИ;
  \item \textbf{T (Technological)}: быстрое устаревание LLM;
    риск лицензионных изменений в~моделях с~открытыми весами;
  \item \textbf{L (Legal)}: 152-ФЗ~\cite{fz_152_2006} (персональные данные);
    авторское право на~ИИ-сгенерированный контент (ГК~РФ~ч.~IV);
  \item \textbf{E (Environmental)}: энергопотребление GPU-кластера;
    углеродный след при~масштабировании.
\end{itemize}

\subsection{Оценка рисков}

\begin{table}[h!]
  \caption{Матрица рисков ИИ-методиста}
  \label{tab:risks}
  \centering
  \begin{tabular}{|p{3.5cm}|p{2cm}|p{2.5cm}|p{4.5cm}|}
  \hline
  \textbf{Риск}
    & \textbf{Вероятность}
    & \textbf{Влияние}
    & \textbf{Митигация} \\
  \hline
  Деградация качества новой версии LLM
    & Средняя
    & Высокое
    & Регрессионный Ragas-тест при~каждом обновлении модели \\
  \hline
  Недоступность GPU-узла
    & Низкая
    & Высокое
    & circuit breaker + резервный переход на~OpenAI~API \\
  \hline
  Утечка данных через LLM-провайдера
    & Низкая
    & Критическое
    & локальное развёртывание; запрет облачных провайдеров для~корпоративного сегмента \\
  \hline
  Неприятие разработчиками курсов
    & Средняя
    & Среднее
    & ручное ревью в~контуре; прозрачность использования ИИ в~интерфейсе \\
  \hline
  Изменение лицензии Qwen / retrieval backend
    & Низкая
    & Высокое
    & LLM~Gateway с~провайдер-агностичным API; альтернативы: Llama~3.3 \\
  \hline
  \end{tabular}
\end{table}

Визуализация рисков на матрице вероятность/влияние приведена на рисунке~\ref{fig:risk-matrix}.

% Матрица рисков ИИ-методиста
\begin{figure}[htbp]
  \centering
  \begin{tikzpicture}
    \begin{axis}[
      width=0.88\textwidth,
      height=0.66\textwidth,
      xmin=0.5, xmax=3.5,
      ymin=0.5, ymax=3.5,
      xtick={1,2,3},
      ytick={1,2,3},
      xticklabels={Низкая, Средняя, Высокая},
      yticklabels={Среднее, Высокое, Критическое},
      xlabel={Вероятность},
      ylabel={Влияние},
      label style={font=\small},
      tick label style={font=\small},
      grid=major,
      grid style={thin, gray!50},
      axis on top,
      clip=false,
    ]
    % Цветовые зоны фона
    % Зелёная зона (низкая вер. + среднее влияние)
    \fill[green!15] (axis cs:0.5,0.5) rectangle (axis cs:1.5,1.5);
    % Жёлтая зона
    \fill[yellow!15] (axis cs:1.5,0.5) rectangle (axis cs:2.5,1.5);
    \fill[yellow!15] (axis cs:0.5,1.5) rectangle (axis cs:1.5,2.5);
    \fill[yellow!15] (axis cs:1.5,1.5) rectangle (axis cs:2.5,2.5);
    % Оранжевая зона
    \fill[orange!15] (axis cs:2.5,0.5) rectangle (axis cs:3.5,1.5);
    \fill[orange!15] (axis cs:2.5,1.5) rectangle (axis cs:3.5,2.5);
    \fill[orange!15] (axis cs:0.5,2.5) rectangle (axis cs:1.5,3.5);
    % Красная зона
    \fill[red!15] (axis cs:1.5,2.5) rectangle (axis cs:2.5,3.5);
    \fill[red!15] (axis cs:2.5,2.5) rectangle (axis cs:3.5,3.5);

    % Риски из матрицы презентации. Небольшой сдвиг по X сохраняет читаемость
    % подписей при одинаковой средней вероятности.
    \addplot[only marks, mark=*, mark size=5pt, red!70] coordinates {(1.82, 3.00)};
    \node[font=\scriptsize, align=center, text width=2.0cm, anchor=south, yshift=4pt]
      at (axis cs:1.82,3.00) {152-ФЗ};

    \addplot[only marks, mark=*, mark size=5pt, orange!80!black] coordinates {(2.18, 2.80)};
    \node[font=\scriptsize, align=center, text width=2.2cm, anchor=south west, xshift=2pt, yshift=2pt]
      at (axis cs:2.18,2.80) {Финансовый\\риск};

    \addplot[only marks, mark=*, mark size=5pt, blue!70] coordinates {(1.82, 2.30)};
    \node[font=\scriptsize, align=center, text width=2.0cm, anchor=north east, xshift=-2pt, yshift=-3pt]
      at (axis cs:1.82,2.30) {ФГОС};

    \addplot[only marks, mark=*, mark size=5pt, purple!70] coordinates {(2.20, 2.15)};
    \node[font=\scriptsize, align=center, text width=2.4cm, anchor=south west, xshift=2pt, yshift=2pt]
      at (axis cs:2.20,2.15) {Авторские\\права};

    \addplot[only marks, mark=*, mark size=5pt, teal!70!black] coordinates {(1.84, 1.55)};
    \node[font=\scriptsize, align=center, text width=2.4cm, anchor=north east, xshift=-2pt, yshift=-3pt]
      at (axis cs:1.84,1.55) {Конкурентный\\риск};

    \addplot[only marks, mark=*, mark size=5pt, green!50!black] coordinates {(2.18, 1.20)};
    \node[font=\scriptsize, align=center, text width=2.0cm, anchor=south west, xshift=2pt, yshift=2pt]
      at (axis cs:2.18,1.20) {Командный\\риск};

    \end{axis}
  \end{tikzpicture}
  \caption{Матрица рисков ИИ-методиста}
  \label{fig:risk-matrix}
\end{figure}


\subsection{Дорожная карта}

\textbf{6 месяцев}: стабилизация оценочного цикла B4, завершение факт-чекинга
и~RAG-оценки, проведение первых решенческих интервью и~пилотного прогона
на~реальных документах заказчиков.

\textbf{12 месяцев}: переход от~чернового MVP к~производственно готовой
версии, поддержка трёх методологий (ADDIE / SAM / Backward Design),
первые платные внедрения и~структурированная пользовательская аналитика.

\textbf{24 месяца}: API для~внешней интеграции, мультиязычность,
масштабирование на~смежные рынки корпоративного обучения.

Дорожная карта представлена на рисунке~\ref{fig:roadmap}.

% Дорожная карта развития ИИ-методиста
\begin{figure}[htbp]
  \centering
  \begin{tikzpicture}[
    phase/.style={draw, rounded corners=3pt, minimum height=1.2cm, text width=3.8cm, align=center, font=\small\bfseries},
    items/.style={font=\footnotesize, align=left, text width=4cm},
  ]
    % Временная ось
    \draw[->, thick] (0, 0) -- (14.5, 0);
    \node[font=\footnotesize, anchor=north] at (0, -0.1) {0};
    \node[font=\footnotesize, anchor=north] at (4.2, -0.1) {6 мес.};
    \node[font=\footnotesize, anchor=north] at (8.4, -0.1) {12 мес.};
    \node[font=\footnotesize, anchor=north] at (14, -0.1) {24 мес.};

    % Фаза 1
    \node[phase, fill=blue!20, minimum width=4cm] (p1) at (2.1, 1.5) {MVP + валидация};
    \draw[thick, blue!50] (0.1, 0.1) -- (0.1, 0.8) -- (p1.south west);
    \draw[thick, blue!50] (4.1, 0.1) -- (4.1, 0.8) -- (p1.south east);
    \node[items, anchor=north west] at (0.1, -0.6) {%
      \textbullet~Стабилизация B4\\
      \textbullet~Проверка фактов\\
      \textbullet~Решенческие интервью%
    };

    % Фаза 2
    \node[phase, fill=green!20, minimum width=4cm] (p2) at (6.3, 1.5) {Промышленная версия 1};
    \draw[thick, green!50] (4.3, 0.1) -- (4.3, 0.8) -- (p2.south west);
    \draw[thick, green!50] (8.3, 0.1) -- (8.3, 0.8) -- (p2.south east);
    \node[items, anchor=north west] at (4.3, -0.6) {%
      \textbullet~3 методологии\\
      \textbullet~Платные внедрения\\
      \textbullet~Аналитика%
    };

    % Фаза 3
    \node[phase, fill=orange!20, minimum width=5.4cm] (p3) at (11.2, 1.5) {Масштабирование};
    \draw[thick, orange!50] (8.5, 0.1) -- (8.5, 0.8) -- (p3.south west);
    \draw[thick, orange!50] (13.9, 0.1) -- (13.9, 0.8) -- (p3.south east);
    \node[items, anchor=north west] at (8.5, -0.6) {%
      \textbullet~Открытый API\\
      \textbullet~Мультиязычность\\
      \textbullet~Смежные рынки%
    };

    % Подпись оси
    \node[font=\small, anchor=north] at (7.25, -2.0) {Время};
  \end{tikzpicture}
  \caption{Дорожная карта развития ИИ-методиста}
  \label{fig:roadmap}
\end{figure}


\subsection{Интеллектуальная собственность}

Программный комплекс подлежит регистрации как~программа
для~ЭВМ (ст.~1261~ГК~РФ~\cite{gk_rf_part4}).
Авторские права на~сгенерированные учебные материалы принадлежат
оператору платформы (ст.~1295~ГК~РФ, служебное произведение).

% ─────────────────────────────────────────────────────────────────
\section{Степень готовности и обратная связь от рынка}
% ─────────────────────────────────────────────────────────────────

\subsection{Текущий статус MVP}

Реализованы: контур обработки документов, CourseStructureAgent (ADDIE),
ручное ревью, ContentCreatorAgent с~RAG, FactCheckAgent.

В~разработке: методологии SAM / Backward~Design, мультимедийный слой
(FLUX~+~Silero), механизм авто-актуализации.

\subsection{Рыночные сигналы}

\begin{itemize}
  \item Проведены 20~B2B-интервью (Customer Development);
  \item 70\,\% респондентов называют время разработки главным ограничением;
  \item 50\,\% уже используют подрядчиков или~внешних методистов;
  \item 60\,\% готовы рассмотреть ИИ-инструмент такого класса;
  \item решенческие интервью, этап писем о~намерениях и~активные пилоты относятся
    к~следующей волне коммерческой валидации.
\end{itemize}

\subsection{Соответствие решения проблеме}

Корректнее говорить не~о~подтверждённом соответствии решения проблеме,
а~о~сильном сигнале наличия проблемы и~раннем интересе к~решению:
проблема валидирована через CustDev, 60\,\% респондентов готовы рассмотреть
ИИ-инструмент. Полноценное подтверждение соответствия
зафиксируется после решенческих интервью на~MVP и~пилотного
использования на~реальных документах.

\subsection{Качественная обратная связь}

Уже на~уровне проблемных интервью выявлены три устойчивых сигнала:
\begin{itemize}
  \item пользователям важнее быстро получить \textit{первый черновик}
    курса, чем сразу идеальный финальный контент;
  \item для~регламентов и~сценариев нормативного обучения критична привязка к~источнику
    и~возможность проверки фактов;
  \item для~B2B-сегмента существенны суверенитет данных, локальное
    развёртывание и~контроль с~участием человека.
\end{itemize}
